
\begin{abstract}

% The text of the abstract goes here.  If you need to use a \section
% command you will need to use \section*, \subsection*, etc. so that
% you don't get any numbering.  You probably won't be using any of
% these commands in the abstract anyway.
% \lipsum[2]
This work develops computational methods for the rapid and accurate assessment of infrastructure following natural hazards.

% The objective of this work is to implement a system for rapid and accurate assessment of infrastructure in the wake of natural hazards.

% This work presents contributions to the field of computational mechanics that target the rapid simulation of structural systems for the purpose of infrastructure assessment in the wake of natural hazards. 

The California Department of Transportation maintains over thirteen thousand bridges, many instrumented with accelerometers that record ground and structural motions during seismic events.
These instruments produce data within seconds of an earthquake, yet the path from recorded accelerations to an actionable assessment of structural condition remains difficult to traverse at the speed and scale that post-earthquake decision-making demands. 

A mean toward this end is to transform sensor data into quantitative indicators of structural condition that can inform decisions about closures, inspections, and repairs.
Two broad strategies are available.
The first, which may be termed forward or physics-based, constructs a mathematical model of the structure and simulates its response to the recorded input motions; predicted response quantities are then compared against capacity thresholds to assess damage.
The second, which may be termed inverse or data-driven, uses recorded input and output motions to identify parameters of a dynamic system; changes in these parameters over time, or deviations from baseline values, serve as indicators of damage.
Both strategies have been explored extensively in the literature, and both have demonstrated value in controlled studies. 
Their integration into operational platforms capable of assessing many structures in near-real-time, however, has proceeded slowly.

A central obstacle is computational. 
Physics-based assessment requires finite element models capable of representing the inelastic behavior that governs damage progression in structural members.
For bridges, these members are predominantly beams and columns, and their accurate simulation under seismic loading involves geometric nonlinearity, cross-sectional warping, and inelastic material response.
There is a wealth of literature addressing the simulation of these phenomena, but many issues remain unresolved. 


Part~I of this dissertation develops a family of beam finite elements that address these issues. 
The central contribution is an enhanced section constitutive framework that improves reconstruction of the three-dimensional strain field while preserving the standard one-dimensional equilibrium equations, with correction amplitudes resolved locally at each section and condensed from the global system.
This constitutive framework is developed in the general setting of geometrically exact beam theory with an arbitrary number of warping modes, and remains valid for linearizations of that model including classical Timoshenko beam theory.
It is then specialized in two directions.
For uniform warping, a trace-based construction is used as a means to derive enhanced shapes that exactly recover Saint~Venant continuum solutions in the elastic regime.
For nonuniform warping, the framework yields enhanced shapes that rectify the torsional equilibrium violation known to arise in standard torsion-warping models.
For inelastic analysis, the framework permits warping effects to be condensed at the section level, so that standard six-degree-of-freedom beam elements can be employed without modification to the element formulation.
The resulting formulations are implemented in \Xara{}, a simulation environment extending the work by \cite{mckenna1997object} that supports automatic differentiation through a hybrid adjoint approach.


Part~II develops a framework for structural health monitoring that motivated and employs the developments of Part~I. 
This framework was implemented for the Bridge Rapid Assessment Center for Extreme Events (\BRACE{}) project with funding from Caltrans to deliver post-earthquake assessments for instrumented highway bridges.
Its computational backend, the Infrastructure Resilience Engine (IRiE), organizes health monitoring around five abstractions: \emph{Assets} representing physical structures, \emph{Events} representing earthquakes or other triggering conditions, \emph{Predictors} implementing physics-based or data-driven analysis methods, \emph{Metrics} quantifying structural condition, and \emph{Evaluations} assembling metrics into actionable reports. 
This architecture permits physics-based simulation and data driven methods to coexist within a unified framework, with results presented through a common interface.
The developments of Part~I contribute the simulation engine for Type~I predictors on this platform.
Following a pilot deployment covering 22 bridges, Caltrans has funded continued development to extend coverage across the California highway network.




% The Bridge Rapid Assessment Center for Extreme Events (\BRACE{}) is a platform developed in collaboration with the California Department of Transportation to provide real-time post-earthquake assessments for instrumented bridges.
% Its computational backend, the Infrastructure Resilience Engine (IRiE), requires finite element models that can be assembled from inventory data, executed rapidly, and embedded in inference workflows.
% The bridges in the current deployment are composed predominantly of beam and column members whose inelastic behavior governs damage progression.
% The beam formulations developed here serve as the primary simulation engine, and the modular structure of the framework---permitting, for instance, a transition from a coarse linear model to a refined inelastic model by changing a single component---directly supports the operational requirements of the platform.



% \section*{Outline}

% The dissertation is organized in two parts.

% Part~I develops finite element formulations for structural members.
% \Cref{sec:Part-I} motivates the theoretical developments and surveys relevant literature. 
% \Cref{sec:genwarp} derives governing equations for a a family of beam theories incorporating warping and finite deformation. 
% % \Cref{sec:transform} presents a framework for geometric transformations that rectifies issues pertaining to strain objectivity that arise in the treatment of finite rotations.
% \Cref{sec:gensec} develops the section-level constitutive framework relating the one-dimensional beam model to the three-dimensional continuum. 

% \Cref{sec:timo} investigates the optimal selection of enhanced shapes in the absence of auxiliary warping fields. The objective is to achieve exact representation of Saint~Venant's continuum solutions through standard beam finite elements. 
% %This precipitates in a novel family of cross section models that exactly represent Saint-Venant's continuum solutions in the linear isotropic regime. 
% % \Cref{sec:library} summarizes the developments and describes their implementation in the \Xara{} and FEDEASLab simulation environments. 
% % \Cref{sec:verify-elastic} presents verification case studies.
% % \Cref{sec:painter} applies the formulations to a Caltrans bridge model.

% Part~II addresses structural health monitoring.
% \Cref{ch:xara} describes the architecture of \Xara{}, with attention to features that support differentiable simulation.
% \Cref{sec:irie} presents the implementation of a structural health monitoring framework for the \BRACE{} project.

\end{abstract}
